\documentclass[aspectratio=169]{beamer}

\usepackage[spanish,es-tabla]{babel}
\usepackage[utf8]{inputenc}
\usepackage[T1]{fontenc}
\usepackage{amsmath}
\usepackage{siunitx}
\usepackage{tikz}
\usetikzlibrary{arrows.meta,babel,positioning}

\usetheme{Madrid}
\usecolortheme{seahorse}
\setbeamertemplate{navigation symbols}{}
\setbeamertemplate{footline}[frame number]
\sisetup{per-mode=symbol,detect-all}

\title{Metasuperficie electromagnética reconfigurable}
\subtitle{Caracterización conceptual en microondas mediante parámetros $S$}
\author{Propuesta experimental conceptual de laboratorio}
\date{\today}

\begin{document}

\begin{frame}
  \titlepage
  \vspace{-0.45cm}
  \begin{block}{Pregunta experimental}
    ¿Cómo se diseñaría la medición de $S_{11}$ y $S_{21}$ para estimar reflexión, transmisión y absorción resonante de una metasuperficie reconfigurable?
  \end{block}
\end{frame}

\begin{frame}{Alcance de la propuesta}
  \Large
  Esta es una \textbf{propuesta experimental conceptual}.

  \vfill
  \begin{block}{Alcance}
    \centering
    \textbf{NO se ejecutará durante el curso.}
  \end{block}

  \vfill
  \begin{itemize}
    \item No se reportan datos medidos.
    \item Se presenta metodología, teoría, datos esperados e incertidumbres.
    \item El valor está en traducir electrodinámica clásica a observables medibles.
  \end{itemize}
\end{frame}

\begin{frame}{Motivación y aplicaciones}
  \Large
  Una metasuperficie delgada puede controlar amplitud, fase, polarización, dirección o absorción de ondas.

  \vfill
  \begin{columns}[T]
    \begin{column}{0.5\textwidth}
      \begin{itemize}
        \item absorbedores electromagnéticos;
        \item shielding y compatibilidad electromagnética;
        \item telecomunicaciones y superficies inteligentes;
        \item control de haz.
      \end{itemize}
    \end{column}
    \begin{column}{0.48\textwidth}
      \begin{itemize}
        \item sensores por corrimiento de resonancia;
        \item radar;
        \item dispositivos THz;
        \item control de polarización.
      \end{itemize}
    \end{column}
  \end{columns}
\end{frame}

\begin{frame}{Objetivos}
  \Large
  \begin{itemize}
    \item Proponer y diseñar un experimento de caracterización electromagnética.
    \item Definir $S_{11}$ y $S_{21}$ como observables principales.
    \item Derivar $R$, $T$, $A$, $f_0$ y $Q$ desde datos esperados.
    \item Formular un procedimiento reproducible, aunque no ejecutado.
    \item Proponer incertidumbres y discutir aplicaciones tecnológicas.
  \end{itemize}
\end{frame}

\begin{frame}{Fundamento: parámetros $S$, $R$, $T$ y $A$}
  \Large
  \begin{align*}
    S_{11}(f)&=\text{coeficiente complejo de reflexión},\\
    S_{21}(f)&=\text{coeficiente complejo de transmisión}.
  \end{align*}

  \vfill
  \[
    R(f)=|S_{11}|^2,
    \qquad
    T(f)=|S_{21}|^2,
    \qquad
    A(f)=1-R(f)-T(f).
  \]

  \vfill
  \normalsize
  El balance energético permite estimar absorción cuando difracción, pérdidas laterales y reflexiones del entorno están controladas.
\end{frame}

\begin{frame}{Impedancia efectiva y condición de absorción}
  \Large
  Para incidencia normal, la reflexión se minimiza cuando
  \[
    Z_\mathrm{eff}\approx Z_0\approx\SI{377}{\ohm}.
  \]

  \vfill
  Con pérdidas internas, la energía no reflejada ni transmitida se disipa:
  \[
    A\ \text{máxima}\quad \Longleftrightarrow \quad R\ \text{baja},\ T\ \text{baja}.
  \]

  \vfill
  \begin{center}
    Adaptación de impedancia + pérdidas controladas = absorción resonante.
  \end{center}
\end{frame}

\begin{frame}{Resonancia LC y modelo dispersivo}
  \Large
  Las celdas de la metasuperficie se modelarían como resonadores sublongitud de onda:
  \[
    f_0=\frac{1}{2\pi\sqrt{LC}},
    \qquad
    Q=\frac{f_0}{\Delta f}.
  \]

  \vfill
  \begin{itemize}
    \item La geometría aporta $L$ y $C$ efectivos.
    \item Las pérdidas conductivas y dieléctricas limitan $Q$.
    \item La reconfiguración modificaría $C$, $f_0$ o la profundidad de absorción.
  \end{itemize}
\end{frame}

\begin{frame}{Montaje propuesto}
  \begin{columns}[T]
    \begin{column}{0.45\textwidth}
      \Large
      \begin{itemize}
        \item VNA de dos puertos.
        \item Dos antenas de microondas.
        \item Metasuperficie PCB.
        \item Absorbentes y control de alineación.
      \end{itemize}
    \end{column}
    \begin{column}{0.53\textwidth}
      \centering
      \begin{tikzpicture}[>=Stealth,scale=0.75]
        \node[draw,rounded corners,minimum width=1.6cm,minimum height=0.8cm] (vna) at (0,0) {VNA};
        \node[draw,minimum width=1.0cm,minimum height=0.8cm,right=1.1cm of vna] (tx) {Ant. 1};
        \node[draw,minimum width=0.2cm,minimum height=2.1cm,right=1.35cm of tx,fill=blue!12] (ms) {};
        \node[above=0.05cm of ms] {MS};
        \node[draw,minimum width=1.0cm,minimum height=0.8cm,right=1.35cm of ms] (rx) {Ant. 2};
        \draw[->,thick] (vna.east) -- (tx.west);
        \draw[->,thick] (tx.east) -- node[above] {incidente} (ms.west);
        \draw[->,thick] (ms.east) -- node[above] {$S_{21}$} (rx.west);
        \draw[<-,thick,red!70!black] (tx.east) -- ++(0.5,-0.7) node[right] {$S_{11}$};
        \draw[->,thick] (vna.south) .. controls (2.5,-1.5) and (4.9,-1.5) .. (rx.south);
      \end{tikzpicture}
    \end{column}
  \end{columns}
\end{frame}

\begin{frame}{Datos esperados, incertidumbres y errores}
  \Large
  Datos esperados:
  \[
    S_{11}(f),\ S_{21}(f),\ R(f),\ T(f),\ A(f),\ f_0,\ Q.
  \]

  \vfill
  \begin{itemize}
    \item Gráficas previstas: módulos, fases, absorción y corrimiento de resonancia.
    \item Incertidumbres: calibración del VNA, cables, alineación, polarización y distancia.
    \item Errores sistemáticos: reflexiones múltiples, difracción de bordes y entorno no anecoico.
  \end{itemize}
\end{frame}

\begin{frame}{Seguridad y conclusiones}
  \Large
  Seguridad propuesta:
  \begin{itemize}
    \item baja potencia de microondas;
    \item conectores RF sin señal aplicada al manipular;
    \item límites de sesgo para elementos reconfigurables;
    \item registro de geometría, potencia y frecuencia.
  \end{itemize}

  \vfill
  \begin{block}{Conclusión conceptual}
    La propuesta muestra cómo condiciones de frontera, impedancia, resonancia, pérdidas y scattering se convierten en magnitudes medibles: $S_{11}$, $S_{21}$, $R$, $T$, $A$, $f_0$ y $Q$.
  \end{block}
\end{frame}

\end{document}
